\documentclass[a4paper,12pt]{article}

\usepackage[T2A]{fontenc}
\usepackage[utf8]{inputenc}
\usepackage[russian]{babel}

\usepackage{amsmath, amssymb, amsthm, mathtools}
\usepackage{helvet}
\renewcommand{\familydefault}{\sfdefault}
\usepackage{icomma}
\usepackage{geometry}
\usepackage{microtype}
\usepackage{tikz}
\usepackage{tkz-euclide}
\usepackage{pgfplots}
\pgfplotsset{compat=1.18}
\usepackage{tabularx, booktabs, longtable}
\usepackage{enumitem}
\usepackage{hyperref}
\usepackage{parskip}
\usepackage{fancyhdr}
\usepackage{setspace}
\usepackage{xcolor}

\definecolor{accent}{HTML}{008080}
\definecolor{darkaccent}{HTML}{004D4D}
\definecolor{textgray}{HTML}{333333}

\geometry{left=2.5cm, right=2.5cm, top=2.5cm, bottom=2.5cm}
\onehalfspacing

\begin{document}

% Титульный лист
\begin{titlepage}
\centering
\vspace*{2cm}
{\Large\bfseries\color{darkaccent} Чек-лист}\par
\vspace{0.5cm}
{{large Тема: Десятичные дроби и текстовые задачи}}\par
\vspace{2cm}
{\large Лёвин Артём Александрович}\par
{\large эксклюзивно для Нади}\par
\vspace{2cm}
{\today}\par
\vfill
\begin{tikzpicture}
\draw[color=accent, line width=1.5pt] 
(0,0) .. controls (2,1.5) and (4,-0.5) .. (6,0.5) .. controls (8,1.5) and (10,-0.5) .. (12,0);
\end{tikzpicture}
\vspace{1cm}
\end{titlepage}

\section*{Чек-лист: Десятичные дроби и текстовые задачи}

\subsection*{Основные правила работы с десятичными дробями}
\begin{itemize}[leftmargin=*]
\item \textbf{Сложение/вычитание}: запятая под запятой.
\item \textbf{Умножение}: перемножить как целые числа, отделить запятой сумму цифр после запятой обоих множителей.
\item \textbf{Деление}: перенести запятую в делимом и делителе на одинаковое число знаков (чтобы делитель стал целым).
\item \textbf{Перевод обыкновенной дроби в десятичную}: разделить числитель на знаменатель.
\end{itemize}

\subsection*{Решение текстовых задач}
\begin{enumerate}[leftmargin=*]
\item Ввести переменную $x$ (обычно меньшую величину).
\item Выразить остальные величины через $x$.
\item Составить уравнение по условию.
\item Решить уравнение, перенося $x$ в одну сторону.
\item Проверить ответ по смыслу задачи.
\end{enumerate}

\subsection*{Типичные ошибки}
\begin{itemize}[leftmargin=*]
\item[$\times$] Неправильно переносят запятую при делении.
\item[$\times$] Забывают посчитать количество знаков после запятой.
\item[$\times$] При решении уравнений не переносят $x$ в одну сторону.
\item[$\times$] Путают произведение и частное при нахождении неизвестного множителя.
\item[$\times$] Не проверяют ответ по смыслу задачи.
\end{itemize}

\newpage
\section*{Тренировочный блок}

\textbf{Уровень 1 (базовый)}

1. Выполните деление: $0{,}783 : 0{,}09$.

2. Вычислите: $2{,}5 + 3{,}75 - 1{,}25$.

3. Переведите в десятичную дробь: $\frac{3}{8}$.

\medskip
\textbf{Уровень 2 (средний)}

4. Решите уравнение: $\frac{0{,}21}{2{,}3x - 41{,}7} = 0{,}47$.

5. Найдите значение: $(12{,}5 \cdot 0{,}8) : 2{,}5$.

6. Представьте в виде десятичной дроби: $\frac{7}{40}$.

\medskip
\textbf{Уровень 3 (выше среднего)}

7. Прыжок в длину в 4{,}2 раза больше прыжка в высоту и на 4{,}48 м больше. Найдите высоту прыжка.

8. Катер прошёл против течения 66{,}15 км за 4{,}9 ч. Собственная скорость 16{,}2 км/ч. Сколько времени понадобится на обратный путь?

9. Вычислите: $\frac{5{,}04 - 0{,}42 \cdot 6}{2{,}1}$.

\medskip
\textbf{Уровень 4 (сложный)}

10. Два велосипедиста выехали одновременно навстречу друг другу из пунктов A и B. Первый едет со скоростью 12{,}5 км/ч, второй --- 14{,}75 км/ч. Через 2{,}4 часа они встретились. Найдите расстояние между пунктами.


\end{document}
